\section{Problem and Claims}
\label{sec:problem}

We study three nested questions:

\begin{enumerate}
  \item \textbf{Model suitability.} Conditioned on an explicit concept
  prefix (\emph{``Personify someone who is terrified of \textit{<concept>}.''}),
  does \mixtral follow the prefix at all on this prompt family? This is a
  prerequisite for any steering claim and is studied separately as a
  full-prefix model comparison against \llama~and \olmoe.
  \item \textbf{Steering transfer.} When the concept prefix is removed from
  the test prompt and the only intervention is the \steermoe routing bias
  learned from the prefix-conditioned routing traces, does the Mixtral output
  move toward the target concept?
  \item \textbf{Steering knob sensitivity.} How does the answer to (2)
  vary as we change the steering coefficient and the number of selected
  positive/negative experts?
\end{enumerate}

\paragraph{Claims.}
We make three narrow claims, all empirical, all about \steermoe-on-\mixtral on
this dataset, and all conservative:

\begin{itemize}
  \item \textbf{Claim 1.} \mixtral~follows explicit fear prefixes
  qualitatively well, comparably to \llama~and visibly better than the
  smaller \olmoe baseline tested on the same prompts. This validates the
  prerequisite and is presented as qualitative evidence (Section~\ref{sec:results}).
  \item \textbf{Claim 2.} Under the question-only test contract, Mixtral +
  \steermoe~produces small, mostly-defocusing shifts relative to the Mixtral
  baseline. The concept-token hit rate moves by at most one case in twenty-five
  across all configurations we tested, while mean response length drops by
  $-2.16$ to $-5.52$ words as the steering coefficient grows.
  \item \textbf{Claim 3.} The signal that drives the \steermoe~plan is
  diffuse on \mixtral: the top twenty $(\text{layer}, \text{expert})$ cells
  account for less than ten percent of the total $|\text{risk diff}|$ mass.
  This is consistent with the question-only behavior and suggests the
  bottleneck is the readout (which tokens the router-logit comparison should
  weight most heavily) rather than the intervention (the additive router bias).
\end{itemize}

\paragraph{What this paper does \emph{not} claim.}
We do not claim \steermoe~fails in general. The original paper's demonstrated
checkpoints differ from \mixtral~in size, instruction tuning, and routing
density, and the upstream evaluation does not use this fear-concept
question-only contract. We also do not claim that any other steering method
would do better on this data; that comparison is the natural follow-up
described in Section~\ref{sec:conclusion}.
